\chapter{Descripci\'on}\label{chapter:descripcion}

La presente secci\'on describe el relevamiento previsto para comprender a los usuarios del dominio y validar que MiGanado responda a necesidades reales de establecimientos bovinos de cr\'ia para la venta. En esta entrega, la descripci\'on se concentra en el proceso de \emph{User Research}, dado que los fundamentos conceptuales se desarrollan en la Secci\'on~\ref{chapter:antecedentes}, el alcance funcional se define en la Secci\'on~\ref{chapter:introduccion} y las decisiones de arquitectura, tecnolog\'ias y validaci\'on se presentan en la Secci\'on~\ref{chapter:metodologia}.

El relevamiento se plantea como una instancia necesaria para evitar que el prototipo se construya \'unicamente a partir de supuestos del equipo. La participaci\'on de productores, veterinarios, administradores y empleados de campo permitir\'a contrastar problemas detectados, reconocer vocabulario propio del sector, identificar restricciones operativas y ordenar prioridades funcionales.

\section{User Research}

\subsection{Objetivo del relevamiento}

El objetivo del \emph{User Research} es obtener informaci\'on cualitativa y cuantitativa sobre las pr\'acticas actuales de gesti\'on ganadera, los problemas frecuentes en el registro de informaci\'on y las condiciones de adopci\'on de una plataforma digital. Esta informaci\'on ser\'a utilizada para validar requerimientos, ajustar la experiencia de usuario y definir criterios de aceptaci\'on del prototipo funcional.

El relevamiento se orientar\'a a comprender c\'omo se registran actualmente los datos individuales de los animales, qu\'e herramientas se emplean para organizar eventos sanitarios, reproductivos y econ\'omicos, y qu\'e dificultades aparecen con mayor frecuencia en la operatoria diaria. Tambi\'en buscar\'a identificar los perfiles que intervienen en la carga, consulta y validaci\'on de informaci\'on, as\'i como las condiciones de conectividad, tiempo disponible, capacitaci\'on y dispositivos que pueden influir en la adopci\'on del sistema.

La informaci\'on obtenida permitir\'a relacionar las necesidades reales del dominio con las funcionalidades previstas para MiGanado. En particular, servir\'a para evaluar la pertinencia de alertas, reportes, asistencia conversacional, estimaciones mediante inteligencia artificial, proyecciones econ\'omicas y mecanismos de uso sin conexi\'on. De esta forma, el \emph{User Research} funcionar\'a como base para priorizar requerimientos y evitar que el prototipo se limite a una definici\'on te\'orica del problema.

\subsection{User Persona}

La definici\'on de \emph{user persona} queda pendiente hasta contar con evidencia validada del relevamiento. De manera preliminar, se identifican cuatro perfiles de usuario relevantes para orientar el an\'alisis: productor o due\~no, administrador de campo, veterinario y empleado operativo. Estos perfiles no deben interpretarse como personas definitivas, sino como hip\'otesis iniciales que deber\'an validarse con informaci\'on obtenida durante el trabajo de campo.

\begin{table}[t]
	\caption{Roles preliminares del sistema}
	\label{tab:roles}
	\centering
	\begin{tabular}{|p{0.24\textwidth}|p{0.61\textwidth}|}
		\hline
		\textbf{Rol} & \textbf{Responsabilidades previstas} \\
		\hline
		Dueno o productor & Visualizar informacion consolidada, indicadores, escenarios economicos y administracion general. \\
		\hline
		Administrador & Gestionar establecimientos, usuarios, lotes, animales y tareas operativas. \\
		\hline
		Veterinario & Consultar historiales sanitarios y reproductivos, registrar controles, tratamientos, tactos y sangrados. \\
		\hline
		Empleado de campo & Cargar datos operativos, consultar animales y registrar eventos autorizados. \\
		\hline
	\end{tabular}
\end{table}

